% Release version, 25 November 2026. Not compiled on the machine where it was written (no LaTeX there): lint only.
\documentclass[11pt]{article}
\usepackage{amsmath,amssymb,amsthm,fullpage}
\usepackage[hidelinks]{hyperref}
\newtheorem{theorem}{Theorem}[section]
\newtheorem{lemma}[theorem]{Lemma}
\newtheorem{proposition}[theorem]{Proposition}
\newtheorem{corollary}[theorem]{Corollary}
\theoremstyle{definition}\newtheorem{conjecture}[theorem]{Conjecture}
\newtheorem{definition}[theorem]{Definition}
\newtheorem{question}[theorem]{Question}
\theoremstyle{remark}\newtheorem{remark}[theorem]{Remark}
\newcommand{\cm}{\operatorname{cm}}
\newcommand{\eps}{\varepsilon}
\newcommand{\RS}{\mathrm{RS}}

\title{The Duchet--Meyniel bound is asymptotically sharp\\ for graphs with independence number two}
\author{Research driven, orchestrated and reviewed by Jack Ryan.\\
Claude (by Anthropic), GPT-6.2 Astra (by OpenAI) and Gemini 4 Neon (by Google)\\
collaborated in finding and writing the proofs.\\
The main theorems are formalized in Lean 4 (Section~\ref{s:verify}).}
\date{25 November 2026}

\begin{document}
\maketitle

\begin{abstract}
Duchet and Meyniel proved in 1982 that every graph $G$ with independence number $\alpha(G)\le2$ has a complete minor of
order at least $|V(G)|/3$. We show that the constant $1/3$ cannot be improved: for every $\beta>0$ there are arbitrarily
large graphs with $\alpha\le2$ in which every connected matching (a matching whose edges are pairwise joined by an edge) has
fewer than $\beta|V(G)|$ edges, and hence $h(G)<(1/3+\tfrac43\beta)|V(G)|+\tfrac23$.
So $\inf h(G)/|V(G)|=1/3$ over nonempty graphs with $\alpha(G)\le2$. This refutes the weak (linear)
connected-matching conjecture of F\"uredi, Gy\'arf\'as and Simonyi and, equivalently, answers in the negative Seymour's
question whether $h(G)\ge(1/3+\eps)|V(G)|$ for some $\eps>0$. Along a sequence of these graphs with $\beta\to0$,
$\chi(G)/h(G)\to3/2$, which by a theorem of B\"ohme, Kostochka and Thomason is the exact supremum of $\chi/h$ for
$\alpha\le2$, and the ratio constants $\rho$ and $\rho_d$ of Kawarabayashi, Wood and Yip for claw-free graphs both equal
$2/3$. The result is a parametric extension of OpenAI's proof of their counterexample to Hadwiger's conjecture (manuscript
dated 23 September 2026). That proof gives the fixed bound $\cm(G)<|V(G)|/100$, and its authors state that the weak
conjecture stays open. Their proof fixes a marginal cap $M=2^{1000}$, and with $M$ fixed its final graph transfer cannot
reach targets below a fixed positive constant; we check that the proof works for every fixed $M\ge1$. Theorem~\ref{t:main}
for rational $\beta$ together with the bound $h(G)<(1/3+2\beta)|V(G)|+2/3$, the upper half of Corollary~\ref{c:third},
and Corollary~\ref{c:conj} are proved in Lean~4 on top of OpenAI's Lean development, with only the three standard axioms.
The other corollaries are short paper arguments on top of the Lean statements and published theorems.
\end{abstract}

\section{Introduction}
Graphs are finite and simple. For a graph $G$ write $\alpha(G)$, $\omega(G)$, $\chi(G)$, $\chi_f(G)$ and $h(G)$ for the
independence number, clique number, chromatic number, fractional chromatic number and Hadwiger number (the largest $t$ with a
$K_t$ minor). Two disjoint edges \emph{touch} if some edge joins their endpoint sets; a \emph{connected matching} is a matching
whose edges touch pairwise, and $\cm(G)$ is the largest size of one. Two disjoint edges \emph{conflict} if all four cross pairs
are non-adjacent; conflicting edges cannot lie in one connected matching.

\begin{theorem}[Duchet--Meyniel \cite{DM}]\label{t:dm}
Every graph $G$ satisfies $h(G)\ge |V(G)|/(2\alpha(G)-1)$. In particular $h(G)\ge|V(G)|/3$ when $\alpha(G)\le2$.
\end{theorem}

Our main result is the following.

\begin{theorem}\label{t:main}
For every $\beta\in(0,1]$ there are arbitrarily large integers $m$ and $m$-vertex graphs $G$ with
\[
 \alpha(G)\le2\qquad\text{and}\qquad \cm(G)<\beta m .
\]
Equivalently, every matching of $G$ with at least $\beta m$ edges contains two conflicting edges.
\end{theorem}

\begin{corollary}\label{c:third}
$\inf\{h(G)/|V(G)|:\alpha(G)\le2,\ |V(G)|\ge1\}=1/3$. More precisely, for every $\eps>0$ there are arbitrarily large graphs with
$\alpha(G)\le2$ and $h(G)<(1/3+\eps)|V(G)|$; by Theorem~\ref{t:dm} no graph with $\alpha\le2$ has $h<|V(G)|/3$.
\end{corollary}

\begin{corollary}\label{c:conj}
The following statements are false.
\begin{enumerate}
\item[(i)] (Linear-CM; weak form of the conjecture of F\"uredi, Gy\'arf\'as and Simonyi \cite{FGS}.) There is $c_0>0$ such
that $\cm(G)\ge c_0|V(G)|$ for every sufficiently large graph $G$ with $\alpha(G)\le2$.
\item[(ii)] ($\mathrm{HC}_\eps$; Seymour \cite{Sey}.) There is $\eps>0$ such that $h(G)\ge(1/3+\eps)|V(G)|$ for every graph
$G$ with $\alpha(G)\le2$.
\end{enumerate}
\end{corollary}

\begin{corollary}\label{c:ratio}
There are graphs $G_j$ with $\alpha(G_j)\le2$ and $|V(G_j)|\to\infty$ such that $\chi(G_j)/h(G_j)\to3/2$ and
$\chi_f(G_j)/h(G_j)\to3/2$.
\end{corollary}

OpenAI's graphs give $\chi(G)>\frac{75}{52}h(G)-\frac{25}{26}$ \cite[cor:hadwiger]{OAIH}, so $\liminf\chi/h\ge75/52$
along their growing witnesses; Corollary~\ref{c:ratio}, along witnesses with $\beta_j\to0$ and $|V(G_j)|\to\infty$, gives
the limit $3/2$. By a theorem of B\"ohme, Kostochka and Thomason \cite{BKT}, $\chi<\frac32h$ for every nonempty graph
with $\alpha\le2$, so $3/2$ is the exact supremum (Remark~\ref{r:upper}). Theorem~\ref{t:main} for rational $\beta$, the
upper half of Corollary~\ref{c:third} and Corollary~\ref{c:conj} are proved in Lean (Section~\ref{s:verify}). The lower
half of Corollary~\ref{c:third} (Theorem~\ref{t:dm}), Corollary~\ref{c:ratio}, Remark~\ref{r:upper},
Corollary~\ref{c:kwy} and Remark~\ref{r:cs} are paper arguments on top of the Lean statements and the cited theorems.
Section~\ref{s:further} adds further paper arguments: the fractional Hadwiger conjectures of Harvey and Wood fail, a question
of Costa, Luu, Wood and Yip on the triangle-free process has a positive answer, and three remarks checked by one reader.

\medskip\noindent\textbf{Status of the result.} The graphs come from the random construction of OpenAI's counterexample to
Hadwiger's conjecture \cite{OAIH}, with scaled parameters. That paper proves Theorem~\ref{t:main} for the single value
$\beta=1/100$ (and so for every $\beta\ge1/100$). Its proof is long (fourteen sections), it has not been refereed, and the
selected consequence $h(G)<26m/75+2/3<m/2\le\chi(G)$ is formalized in OpenAI's Lean library \cite{OAIL}; the repository
\cite{Bon} ports that Lean proof (with its 288 upstream modules) and reports a kernel replay. The present paper is a parametric extension
of OpenAI's proof. What we add is (a) an audit showing that the marginal cap $M=2^{1000}$ of the raw supersaturation
theorem, the one constant that keeps $\beta$ away from $0$ in \cite{OAIH}, can be replaced by any fixed $M\ge1$
(Proposition~\ref{p:uniform}); (b) the corresponding change of the
graph transfer (Lemma~\ref{l:transfer}); (c) a Lean~4 proof of the main statements, obtained from OpenAI's Lean proof by
the same two changes (Section~\ref{s:verify}). In the Lean proof, OpenAI's formal supersaturation theorem is used
unchanged; it already holds for every value of the cap. So the Lean proof does not depend on the paper audit (a). A second, independent route
through the companion Colin de Verdi\`ere paper \cite{OAIC} is sketched in Section~\ref{s:second}.

\subsection{History}
Hadwiger conjectured in 1943 that $h(G)\ge\chi(G)$ \cite{Had}. For $\alpha(G)\le2$ every colour class has at most two
vertices, so $\chi(G)\ge|V(G)|/2$, and Plummer, Stiebitz and Toft \cite[Theorem~3.1]{PST} showed that the restriction of the
conjecture to $\alpha(G)\le2$ is equivalent to the statement that every such graph has a $K_{\lceil |V(G)|/2\rceil}$ minor.
Theorem~\ref{t:dm} gives $|V(G)|/3$, and for a long time the improvements were of lower order: Fox \cite[Theorem~2]{Fox}
proved $h(G)\ge |V(G)|/3+c|V(G)|^{4/5}(\log|V(G)|)^{1/5}$ for $\alpha(G)\le2$, and his argument also gives connected
matchings of size $\Omega(|V(G)|^{4/5}(\log |V(G)|)^{1/5})$.

F\"uredi, Gy\'arf\'as and Simonyi \cite{FGS} introduced connected matchings in this setting. They conjectured the sharp
statement that every graph with $\alpha=2$ on $4t-1$ vertices has a connected matching of size $t$
\cite[Conjecture~1.5]{FGS} (verified for $t\le17$ by them \cite[Theorem~1.6]{FGS} and for $t\le22$ by Chen and Deng
\cite{CD}; Cambie \cite{Cam} showed that Hadwiger's conjecture implies it), and they stated the weaker linear form: some
constant $c$ gives every graph with $ct$ vertices and $\alpha=2$ a connected matching of size $t$ \cite[Conjecture~1.4]{FGS}.
This is Corollary~\ref{c:conj}(i) (Linear-CM, \cite[Conjecture~1.8]{CLWY}; \cite[Conjecture~7.1]{CFS}). They state
Seymour's problem of improving $1/3$ as \cite[Conjecture~1.3]{FGS}, and they report that Thomass\'e noted, at a 2002
problem session in Odense, that the two conjectures are equivalent, with a proof by Kawarabayashi, Plummer and Toft
\cite{KPT} (see also \cite[Section~7]{CFS} and \cite[Section~1]{CLWY}). The question is Open question~4.8 of Seymour's
survey \cite{Sey}: whether some constant $c>1/3$ gives every graph $G$ with $\alpha(G)=2$ a $K_t$ minor with
$t\ge c|V(G)|$. Costa, Luu, Wood and Yip call it $\mathrm{HC}_\eps$ \cite[Conjecture~1.7]{CLWY}, and Conlon, Fox and Sudakov
write that Seymour and Mader have both asked it \cite[Section~7]{CFS}. Yip \cite[Theorem~7]{Yip} proved a dense variant:
linear-size matchings in which the proportion of non-touching pairs is arbitrarily small. (We have not read \cite{KPT}
itself; its role here follows \cite{FGS,CLWY}.)

In a manuscript dated 23 September 2026 and made public in October 2026, OpenAI gave a construction of graphs with
$\alpha(G)\le2$ and $\cm(G)<|V(G)|/100$ \cite[thm:main]{OAIH},
which disproves Hadwiger's conjecture and the sharp FGS conjecture. The authors write
\cite[subsection ``Historical context'']{OAIH}: ``The fixed coefficient in [the main theorem] leaves the weaker conjecture for
smaller positive coefficients unresolved; it does not assert $\mathrm{cm}(G)=o(|V(G)|)$.'' Theorem~\ref{t:main} supplies
exactly this, by the same construction.

\section{The OpenAI construction as a black box}\label{s:box}
We use \cite{OAIH} only through the following objects and facts; section and label names refer to that paper.
For each $n$ there are a finite set $\Omega_n$ of \emph{raw vertices}, a strictly positive probability measure $\mu$ on
$\Omega_n$, and a symmetric relation of \emph{holes} on $\Omega_n$ with no loops and no triangles (no three raw vertices
pairwise in holes) \cite[lem:hole-triangle-free]{OAIH}. The ambient dimension is $N=M_0n$. A \emph{unit} is an ordered pair
$(x,y)\in\Omega_n^2$ that is not a hole; two units \emph{conflict} if all four cross pairs are holes. For a probability law
$\sigma$ on units let $\sigma_1,\sigma_2$ be its marginals. The construction constants are fixed in the order
\begin{equation}\label{eq:order}
 \begin{gathered}
 (C_0,g,D,M,k_{\max},\zeta,K_1)\to(r,L_0,d_0,u_0,K,\text{tolerances},p_0)\\
 \to(R_*,j_*,B_0,B_*,b,p_*,A,s_0,J)\to(r_0,h)\to M_0\to n
 \end{gathered}
\end{equation}
\cite[Section~14]{OAIH}, with $C_0=1000$, $g=10^9+1$, $D=4C_0g$ and $M=2^{1000}$.

\begin{definition}
For a real $M\ge1$ let $\RS(M)$ be the statement: \emph{one can choose the constants after $M$ in \eqref{eq:order} so that for
all sufficiently large $n$, every law $\sigma$ on units with}
\begin{equation}\label{eq:caps}
 \sigma_1\le M\mu,\qquad\sigma_2\le M\mu,\qquad \sigma\le2^{DN}\mu^{\otimes2}
\end{equation}
\emph{gives two independent $\sigma$-units that conflict with probability at least $2^{-100gN}$.}
\end{definition}

\begin{theorem}[Raw supersaturation {\cite[thm:raw-supersaturation]{OAIH}}]\label{t:rs}
$\RS(2^{1000})$ holds.
\end{theorem}

$\RS(M)$ becomes stronger as $M$ grows, since the class of admissible laws grows. It has a simple necessary condition:
if $S\subseteq\Omega_n$ contains no hole and $\mu(S)\ge1/M$, then $\sigma=(\mu|_S/\mu(S))^{\otimes2}$ satisfies
\eqref{eq:caps} for large $n$ and has no conflicts. Indeed, every pair in $S\times S$ is a unit (holes have no loops and $S$
contains no hole), both marginals equal $\mu|_S/\mu(S)\le M\mu$, the joint law is at most
$\mu^{\otimes2}/\mu(S)^2\le M^2\mu^{\otimes2}\le2^{DN}\mu^{\otimes2}$ once $2^{DN}\ge M^2$, and two units inside
$S\times S$ cannot conflict because their cross pairs lie in $S\times S$ and are not holes. So $\RS(M)$ for every $M$ forces
the $\mu$-mass of hole-free sets to tend to $0$; this matches $\omega(G)\le2\cm(G)+1<2\beta m+1$ for the graphs of
Theorem~\ref{t:main}.

\section{The marginal cap is free}\label{s:uniform}
\begin{proposition}[conditional on \cite{OAIH}]\label{p:uniform}
$\RS(M)$ holds for every real $M\ge1$.
\end{proposition}

\begin{proof}[Proof (audit)]
In \eqref{eq:order} the cap $M$ is in the first group, and none of $k_{\max},\zeta,K_1$ or of the constants of the second
group ($r,L_0,d_0,u_0,K$ are explicit in $D$ and $\zeta$; the tolerances are numbers; $p_0$ depends on $r,D$ and numbers)
depends on $M$. Every later constant may depend on $M$. It therefore suffices that the proof of Theorem~\ref{t:rs} never
needs $M$ to be smaller than a quantity fixed before $M$, and that every use of $M$ is paid by a constant chosen after $M$.
The proof of Theorem~\ref{t:rs} is by contradiction along a sequence of laws \cite[Section~13]{OAIH}, and the marginal cap
enters only in the following places (line numbers of the TeX sources of \cite{OAIH}):
\begin{enumerate}
\item Peeling (\texttt{03-frame-laws}, l.280--330) uses only the joint cap; after a restriction of mass $2^{-o(N)}$ the
marginals are at most $M2^{o(N)}\mu$: a constant factor.
\item Phase estimates (\texttt{07-phases}). The fixed positive-mass restrictions (l.57--66) have mass bounded below in terms
of $D,g,K_1$ and numbers only, so the marginal bound after them is $L\mu$ with $L$ a constant multiple of $M$. Then $L$
multiplies an exceptional mass $2^{-\Omega(N)}$ (l.290--303), is absorbed into $M_0$ (l.380--386), and enters the
requirement $L\,O(r)\,2^{2r-h}<10^{-4}$ (l.395--400), which is met by the choice of $h$ after $M$. In the injection lemma
(l.560--640) it multiplies $2^{-\Omega(N)}$.
\item Preparation (\texttt{08-preparation}, l.158--160, 192--197, 351, 444--447, 497--499, 614): the marginal bound appears as a
fixed factor in front of an $o(1)$ reference error, or as an $M2^{o(N)}$ inflation of $2^{-\Omega(n)}$ errors.
\item Histograms (\texttt{09-histograms}, l.510--518): caps $M2^{a_nN}\mu$; the authors remark there that the value of the
fixed factor $M$ is immaterial. The same mixed bound is used in \texttt{11-common-limit} (l.326) and \texttt{13-rare-overlap}
(l.77, 184), only through the typicality lemma, with errors $2^{-\Omega(n)}$.
\item The collision criterion (\texttt{06-collision}) and the status-span section (\texttt{12-status-span}, l.555--570) use
exponential scales only.
\item In the probability-scale table (\texttt{14-parameters}, l.140--150, 399--425) the one-vertex density factor enters a
lower bound on $h$; all other entries treat $M$ like $2^{o(N)}$.
\end{enumerate}
No step compares $M$ with a constant chosen before $M$. Hence the proof of Theorem~\ref{t:rs} goes through verbatim with any
fixed $M\ge1$ in place of $2^{1000}$, the constants after $M$ being chosen for that $M$.
\end{proof}

\noindent Proposition~\ref{p:uniform} is the step that gives the extra strength. Its checks are listed in
Section~\ref{s:verify}. For the Lean proof it is not needed: see Remark~\ref{r:lean}.

\begin{remark}[machine evidence]\label{r:lean}
In OpenAI's Lean development \cite{OAIL} (module family \texttt{OAI.Combinatorics.HadwigerCounterexample}) the marginal cap
is a free real field \texttt{L} of the structure \texttt{PreparationParameters}; the admissibility predicate asks only
$0\le\texttt{L}$; the phase calendar \texttt{phase\_calendar g r k d D L}, under the hypothesis $0\le\texttt{L}$, concludes a statement eventually in
$h$, so $h$ is chosen after $\texttt{L}$, and its retained mass does not contain \texttt{L}; and
\texttt{eventual\_conflict\_lower\_bound} is proved for arbitrary parameters. The value $2^{1000}$ occurs only in
\texttt{ParameterExistence.lean}, where the equation $\texttt{b.L}=2^{1000}$ is used only to be exported, and in the final
graph transfer. We restated the two parameter-existence lemmas with a parameter in place of $2^{1000}$ and the graph
transfer with the threshold $m/Q$, and the resulting Lean proof of the main statements compiles
(Section~\ref{s:verify}). In that proof the supersaturation theorem \texttt{eventual\_conflict\_lower\_bound} is OpenAI's,
unchanged. So the formal proof covers the statement of Proposition~\ref{p:uniform} in OpenAI's Lean model of the
construction, for every cap $2^{12Q+1}$.
\end{remark}

\section{The scaled graph transfer}\label{s:transfer}
The following two facts are proved in \cite[Section~3]{OAIH}; their proofs do not use the value of $M$.

\begin{lemma}[terminal cut {\cite[lem:terminal-cut]{OAIH}}]\label{l:cut}
If $R\subseteq\Omega_n^2$ supports no probability law satisfying \eqref{eq:caps}, there are $S\subseteq\Omega_n$ and
$E_0\subseteq\Omega_n^2$ with $\mu(S)<1/M$, $\mu^{\otimes2}(E_0)<2^{-DN}$, such that every pair in $R$ has an endpoint
in $S$ or lies in $E_0$.
\end{lemma}

\begin{lemma}[entropy fingerprints {\cite[proof of prop:raw-to-graph]{OAIH}}]\label{l:fp}
Assume the conclusion of $\RS(M)$ at $n$. There is a family $\mathcal T$ of sets of units with
$\log|\mathcal T|=O(N^3\,2^{100gN})$, the constant independent of $M$, such that no $T\in\mathcal T$ supports a law
satisfying \eqref{eq:caps}, and every set of pairwise non-conflicting units is contained in some $T\in\mathcal T$.
\end{lemma}
\noindent (The fingerprint length is $1+\lceil DN\log2/(-\log(1-2^{-100gN}))\rceil$, because the relative entropy of a
feasible law lies in $[0,DN\log2]$; $M$ does not enter.)

\begin{lemma}[scaled transfer]\label{l:transfer}
Let $\beta\in(0,1]$ and $M\ge4e/\beta$. Assume the conclusion of $\RS(M)$ at a sufficiently large $n$. Let $m=2^{C_0gN}$, draw
$x_1,\dots,x_m$ independently from $\mu$, and let $G$ be the graph on $\{1,\dots,m\}$ in which $i\ne j$ are adjacent iff
$(x_i,x_j)$ is not a hole. Then $\alpha(G)\le2$ always, and with probability $1-\exp(-\Omega_\beta(m))$ every matching of
$G$ with at least $\beta m$ edges contains two conflicting edges; in particular $\cm(G)<\beta m$.
\end{lemma}

\begin{proof}
An independent triple of $G$ consists of three positions whose types are pairwise in holes; the types are distinct because
holes have no loops, which contradicts triangle-freeness. So $\alpha(G)\le2$.

Put $k=\lceil\beta m/2\rceil$. For each $T\in\mathcal T$ fix $S_T,E_T$ by Lemma~\ref{l:cut}. For a fixed $T$,
\[
 \Pr\bigl(|\{i:x_i\in S_T\}|\ge k\bigr)\le\binom mk\mu(S_T)^k\le\Bigl(\frac{em}{kM}\Bigr)^k\le\Bigl(\frac{2e}{\beta M}\Bigr)^k\le2^{-k}.
\]
A fixed family of $k$ disjoint ordered pairs of distinct positions consists of $2k$ independent $\mu$-draws, so lies in
$E_T$ with probability less than $2^{-kDN}$; there are at most $m^{2k}$ such families, and
$m^{2k}2^{-kDN}=2^{(2C_0g-D)Nk}=2^{-2C_0gNk}$. By Lemma~\ref{l:fp} and the union bound, the probability that some
$T\in\mathcal T$ has at least $k$ positions in $S_T$ or $k$ disjoint position pairs in $E_T$ is at most
\[
 \exp\bigl(O(N^3 2^{100gN})\bigr)\bigl(2^{-k}+2^{-2C_0gNk}\bigr)=\exp(-\Omega_\beta(m)),
\]
because $k\ge\beta m/2=\beta2^{1000gN}/2$ while $N^32^{100gN}=o(m)$.

On the complementary event, let $F$ be a matching of $G$ with $t\ge\beta m$ edges and assume no two of its edges conflict.
Orient each edge $ij$ of $F$ and let $I$ be the set of units $(x_i,x_j)$ so obtained (repetitions allowed). Two units of $I$
coming from different edges do not conflict, since conflicting types would make the two edges conflict; a unit does not
conflict with itself because holes have no loops. So $I$ lies in some $T\in\mathcal T$, and every edge $ij$ of $F$ has
$x_i\in S_T$ or $x_j\in S_T$ or $(x_i,x_j)\in E_T$. The edges of $F$ are disjoint, so at most $k-1$ of them meet a position
of type in $S_T$, and the others form at most $k-1$ disjoint ordered position pairs in $E_T$. Hence
$t\le2k-2<\beta m$, a contradiction. A connected matching has no conflicting pair, so $\cm(G)<\beta m$.
\end{proof}

\section{Proofs of the main results}\label{s:proofs}
\begin{proof}[Proof of Theorem~\ref{t:main}]
Given $\beta$, put $M=\max\{2^{1000},\lceil4e/\beta\rceil\}$. Proposition~\ref{p:uniform} gives $\RS(M)$, with constants
(in particular $M_0$) depending on $\beta$. For every sufficiently large $n$, Lemma~\ref{l:transfer} gives a graph on
$m=2^{C_0gM_0n}$ vertices with the two properties; $m\to\infty$ with $n$.
\end{proof}

\begin{lemma}[{\cite[prop:matching-minor]{OAIH}}, every graph]\label{l:mm}
$h(G)\le(|V(G)|+4\cm(G)+2)/3$.
\end{lemma}

\begin{proof}[Proof of Corollaries~\ref{c:third} and \ref{c:conj}]
For the graphs of Theorem~\ref{t:main}, Lemma~\ref{l:mm} gives $h(G)<(1+4\beta)m/3+2/3$. With $\beta=\eps/2$ this is below
$(1/3+\eps)m$ once $m>2/\eps$. Theorem~\ref{t:dm} gives the lower bound $1/3$. Corollary~\ref{c:conj}(i) fails at
$\beta<c_0$ and (ii) fails by Corollary~\ref{c:third}; the equivalence of (i) and (ii) \cite{KPT} is not needed.
\end{proof}

\begin{proof}[Proof of Corollary~\ref{c:ratio}]
Let $G$ be as in Theorem~\ref{t:main}, with $m$ vertices. The vertices of a clique can be paired into $\lfloor\omega/2\rfloor$
pairwise touching edges, so $\omega(G)\le2\cm(G)+1<2\beta m+1$. Since $\alpha(G)\le2$, colour classes are single vertices or
non-adjacent pairs, so $\chi(G)=m-\nu(\overline G)$, where $\nu(\overline G)$ is the matching number of the complement. The
vertices missed by a maximum matching of $\overline G$ are pairwise adjacent in $G$, so $m-2\nu(\overline G)\le\omega(G)$ and
\begin{gather*}
 \frac m2\le\chi_f(G)\le\chi(G)\le\frac{m+\omega(G)}2<\Bigl(\frac12+\beta\Bigr)m+\frac12,\\
 \frac m3\le h(G)<\frac{(1+4\beta)m}3+\frac23 .
\end{gather*}
(The lower bound for $\chi_f$ is in the proof of \cite[cor:fractional-hadwiger]{OAIH}.) Take $\beta_j\to0$ and $m_j\to\infty$; both ratios tend to $3/2$.
\end{proof}

\begin{remark}[$3/2$ is the end]\label{r:upper}
B\"ohme, Kostochka and Thomason \cite[Lemma~2.3]{BKT} proved $h(G)\ge(4\chi(G)-|V(G)|)/3$ for every graph with
$\alpha(G)\le2$ (their Theorem~1.1 extends this to all graphs). Since $|V(G)|\le2\chi(G)$ when $\alpha(G)\le2$, this gives
$\chi(G)\le\frac32h(G)$, and with their Lemma~2.1 ($h(G)\ge1+|V(G)|/3$ for $\alpha(G)\le2$ and $|V(G)|\ge5$) the
inequality is strict. Hence, by Corollary~\ref{c:ratio},
\[
 \sup\{\chi(G)/h(G):\alpha(G)\le2,\ |V(G)|\ge1\}=\sup\{\chi_f(G)/h(G):\alpha(G)\le2,\ |V(G)|\ge1\}=3/2,
\]
and neither supremum is attained. The same lemma shows that any family with $h/|V|\to1/3$ must have
$\chi=(1/2+o(1))|V|$; for our graphs this also follows from the bound $\chi\le(m+\omega)/2$ in the proof above.
Explicitly, Lemmas~2.1 and 2.3 of \cite{BKT} give $h\ge1+n/3$ and $h\ge(4\chi-n)/3$; their sum gives
$\chi\le\frac32h-\frac34$.
\end{remark}

\begin{corollary}[claw-free graphs]\label{c:kwy}
Let $\rho$ and $\rho_d$ be the infima of $h(G)/\chi(G)$ and of $h_d(G)/\chi(G)$ over non-empty claw-free graphs $G$, where
$h_d$ is the largest order of a dominating clique model, as in \cite[Section~1.1]{KWY}. Then $\rho=\rho_d=2/3$.
\end{corollary}

\begin{proof}
Kawarabayashi, Wood and Yip record $2/3\le\rho_d\le\rho\le2/3+2/75$ and ask, as their Open Problem~(3), ``What are the
constants $\rho$ and $\rho_d$?'' \cite[Sections~1.1 and~7]{KWY}. A graph with $\alpha(G)\le2$ is claw-free. For the graphs of
Theorem~\ref{t:main}, $\chi(G)\ge m/2$ and $h(G)<(1+4\beta)m/3+2/3$, so $h(G)/\chi(G)<2/3+\tfrac83\beta+4/(3m)$. Hence
$\rho\le2/3$, and so $\rho=\rho_d=2/3$. (The lower bound $2/3$ is \cite[Theorem~5]{KWY}, which generalizes a theorem of
Chudnovsky and Fradkin \cite{CF}.)
\end{proof}

\begin{remark}[the clique threshold $|V|/4$]\label{r:cs}
Chudnovsky and Seymour \cite[Theorem~1.3]{CS} proved that a graph with $\alpha(G)\le2$ has a $K_{\lceil |V(G)|/2\rceil}$
minor when $\omega(G)\ge|V(G)|/4$ ($|V(G)|$ even) or $\omega(G)\ge(|V(G)|+3)/4$ ($|V(G)|$ odd) (statement as given in
\cite[subsection ``Historical context'']{OAIH}). The constant $1/4$ cannot be
lowered. Let $G_0$ be a graph of Theorem~\ref{t:main} with $n_0$ vertices, and let $G$ be the join of $G_0$ and $K_q$ (every
vertex of $G_0$ joined to every vertex of $K_q$), with $n=n_0+q$. Then $\alpha(G)\le2$, $\omega(G)=\omega(G_0)+q$, and
$h(G)\le h(G_0)+q$, since deleting one vertex lowers the Hadwiger number by at most one. So $h(G)<n/2$ whenever
$q<n_0-2h(G_0)$, and $n_0-2h(G_0)>(1/3-\tfrac83\beta)n_0-\tfrac43$. With $q=\lceil(1/3-3\beta)n_0\rceil$ and $n_0$ large,
$G$ has no $K_{\lceil n/2\rceil}$ minor and $\omega(G)/n\ge q/(n_0+q)$, which tends to $1/4$ as $\beta\to0$. (OpenAI's
fixed bound $h<26n_0/75+2/3$ alone gives the same with $\omega/n$ up to about $23/98$.)
\end{remark}

\begin{remark}[larger independence number]
Disjoint unions of copies of our graphs give $h\approx|V(G)|/(3j)$ with $\alpha=2j$, which does not reach the
Duchet--Meyniel value $|V(G)|/(4j-1)$. For $\alpha\ge3$ the bound of Theorem~\ref{t:dm} is improvable: Kawarabayashi and
Song replaced $2\alpha-1$ by $2\alpha-2$ \cite[Theorem~4.4]{Sey} (an earlier improvement to $2\alpha-3/2$ is due to
Kawarabayashi, Plummer and Toft \cite{KPT}); for large $\alpha$ the best factor is $0.51338\,|V(G)|/\alpha(G)$ of
Balogh and Kostochka \cite{BK}. So $\alpha=2$ is the case in which the bound of Duchet and Meyniel is asymptotically
sharp.
\end{remark}

\section{Further consequences}\label{s:further}
The results of this section are paper arguments on top of Theorem~\ref{t:main}. They are not formalized in Lean.

\subsection{The fractional Hadwiger number}
Two subgraphs of $G$ \emph{touch} if they share a vertex or some edge of $G$ joins them. A \emph{bramble} is a set of
connected subgraphs of $G$ that pairwise touch. A weighting $w$ of a bramble $\mathcal B$ with nonnegative reals is
\emph{admissible} if $\sum_{X\in\mathcal B,\ v\in V(X)}w(X)\le1$ for every vertex $v$. The \emph{fractional Hadwiger number}
$h_f(G)$ is the largest total weight $\sum_{X\in\mathcal B}w(X)$ of an admissible weighting of a bramble
\cite[Sections~3 and 11]{HW}; the same parameter was introduced by Fox \cite{Fox11} and by Pedersen (see \cite[Section~11]{HW}).
The branch sets of a $K_{h(G)}$ model with weight $1$ show $h(G)\le h_f(G)$. Harvey and Wood
\cite[Section~12, Conjectures~23 and 24]{HW} conjectured $\chi(G)\le h_f(G)$ and $\chi_f(G)\le h_f(G)$ for every graph $G$.

\begin{proposition}\label{p:hf}
Every graph $G$ on $n$ vertices satisfies $3h_f(G)\le n+4\cm(G)+2$. Equality holds for complete graphs of odd order.
\end{proposition}

\begin{proof}
Fix a bramble $\mathcal B$ with an admissible weighting $w$; we may drop the members of weight $0$. Let $p$, $q$ and $r$ be the
total weights of the members with one vertex, two vertices and at least three vertices. Summing the vertex constraints gives
$\sum_X w(X)|V(X)|\le n$, so $p+2q+3r\le n$ and
\[
 3(p+q+r)\le n+2p+q .
\]
Two one-vertex members $\{u\}$ and $\{v\}$ touch only if $u=v$ or $uv\in E(G)$. So the set $K$ of vertices $v$ with
$\{v\}\in\mathcal B$ is a clique; let $k=|K|$. A two-vertex member is connected, so it is an edge. Let $q_K$ be the total
weight of the edge members that meet $K$. The constraints at the vertices of $K$ give $p+q_K\le k$; since every member weighs
at most $1$, also $p\le k$. Let $M$ be a maximal set of pairwise disjoint edge members that do not meet $K$. Every edge
member that does not meet $K$ meets $V(M)$, so the constraints at the $2|M|$ vertices of $V(M)$ give $q-q_K\le2|M|$. Hence
\[
 2p+q=p+(p+q_K)+(q-q_K)\le 2k+2|M| .
\]
Now pair the vertices of $K$ into $\lfloor k/2\rfloor$ disjoint edges and add $M$. The result is a connected matching. Indeed
two edges inside the clique $K$ touch; two edges of $M$ touch because they are members of $\mathcal B$; and a vertex
$v\in K$ is adjacent to an endpoint of every edge of $M$, because the member $\{v\}$ touches every member and $v$ lies on no
edge of $M$. So $\lfloor k/2\rfloor+|M|\le\cm(G)$, and with $k\le2\lfloor k/2\rfloor+1$,
\[
 2p+q\le 4\lfloor k/2\rfloor+2+2|M|\le4\cm(G)+2 .
\]
Hence $3h(\mathcal B,w)\le n+4\cm(G)+2$, and the inequality follows by maximizing over $(\mathcal B,w)$. For $G=K_n$ with $n$
odd, $h_f(G)=n$ and $\cm(G)=(n-1)/2$, so both sides equal $3n$.
\end{proof}

\begin{corollary}\label{c:hf}
\begin{enumerate}
\item[(i)] For every $\eps>0$ there are arbitrarily large graphs $G$ with $\alpha(G)\le2$ and
$h_f(G)<(1/3+\eps)|V(G)|$.
\item[(ii)] Conjectures~23 and 24 of Harvey and Wood are false. Over nonempty graphs with $\alpha(G)\le2$,
\[
 \sup\chi(G)/h_f(G)=\sup\chi_f(G)/h_f(G)=3/2 ,
\]
and neither supremum is attained.
\end{enumerate}
\end{corollary}

\begin{proof}
(i) For the graphs of Theorem~\ref{t:main}, Proposition~\ref{p:hf} gives $h_f(G)<(1+4\beta)m/3+2/3$; take $\beta=\eps/2$ and
$m>2/\eps$.

(ii) A graph with $\alpha(G)\le2$ and $m$ vertices has $\chi(G)\ge\chi_f(G)\ge m/2$, so both conjectures fail as soon as
$h_f(G)<m/2$. This already happens for OpenAI's graphs: from $\cm(G)<m/100$, Proposition~\ref{p:hf} gives
$h_f(G)<26m/75+2/3<m/2$ for $m\ge5$. For the value of the suprema, take the graphs of Theorem~\ref{t:main} with
$\beta\to0$: by (i) and the proof of Corollary~\ref{c:ratio}, $h_f=(1/3+o(1))m$ and $m/2\le\chi_f\le\chi\le(1/2+o(1))m$, so
both ratios tend to $3/2$. Conversely, $\chi_f(G)\le\chi(G)<\frac32h(G)\le\frac32h_f(G)$ for every nonempty graph with
$\alpha(G)\le2$ (Remark~\ref{r:upper}), so $3/2$ is not attained.
\end{proof}

\begin{remark}
Fox \cite[Section~4]{Fox11} remarks that, for graphs with independence number two, improving the bound $h\ge|V(G)|/3$ to
$(1/3+\eps)|V(G)|$ is equivalent to the same improvement for $h_f$. With Corollary~\ref{c:conj} this already gives the
refutation of Conjectures~23 and 24 in qualitative form. Proposition~\ref{p:hf} gives it with explicit constants, from
OpenAI's fixed bound alone, and the exact value $3/2$ of the suprema needs Theorem~\ref{t:main}.
\end{remark}

\subsection{The triangle-free process}
Costa, Luu, Wood and Yip \cite[Question~4.8]{CLWY} (numbered Question~21 in the HTML version of arXiv:2512.17114v2) ask: let
$G$ be the complement of a random graph generated by the triangle-free process with $|V(G)|\ge4t-1$; as $t\to\infty$, does
$G$ contain a connected matching of size $t$ asymptotically almost surely? The answer is yes. For a graph $F$ write $e(F)$ for
its number of edges, $\Delta(F)$ for its maximum degree and $C_4(F)$ for its number of $4$-cycles (subgraphs isomorphic to
$C_4$).

\begin{proposition}\label{p:pairing}
Let $G$ be a graph on $n\ge3$ vertices, let $F=\overline G$, and let $N$ be the least even integer with $N\ge n$. Then
\[
 \cm(G)\ \ge\ \Bigl\lfloor\frac n2\Bigr\rfloor-\frac{e(F)}{N-1}-\frac{C_4(F)}{(N-1)(N-3)} .
\]
\end{proposition}

\begin{proof}
If $n$ is odd, add one isolated dummy vertex; this gives $N$ vertices. Choose a uniformly random perfect matching $P$ of
these $N$ vertices and remove the pair that contains the dummy, if any. The remaining $\lfloor n/2\rfloor$ pairs are
disjoint pairs of vertices of $G$. A fixed pair $\{u,v\}$ is in $P$ with probability $1/(N-1)$, so the expected number of
pairs that are edges of $F$ is $e(F)/(N-1)$. Two disjoint pairs $\{a,b\}$ and $\{c,d\}$ that are edges of $G$ fail to touch
exactly when $ac$, $ad$, $bc$ and $bd$ are all edges of $F$, that is, when $a,c,b,d$ is a $4$-cycle of $F$ whose two pairs of
opposite vertices are $\{a,b\}$ and $\{c,d\}$. Each $4$-cycle of $F$ determines its two pairs of opposite vertices, and both
are in $P$ with probability $1/((N-1)(N-3))$. So the expected number of such failures is at most $C_4(F)/((N-1)(N-3))$.
Delete from $P$ every pair that is an edge of $F$, and then one pair from every remaining failing couple. The rest is a
connected matching of $G$, and its expected size is at least the right-hand side. So some choice of $P$ does at least as
well.
\end{proof}

\begin{corollary}\label{c:tfp}
Let $F$ be the final graph of the triangle-free process on $n$ vertices and $G=\overline F$. Then asymptotically almost
surely $\cm(G)\ge n/2-O(\sqrt n\,\log^{3/2}n)$. In particular the answer to \cite[Question~4.8]{CLWY} is yes.
\end{corollary}

\begin{proof}
$F$ is triangle-free, so every neighbourhood in $F$ is an independent set and $\Delta(F)\le\alpha(F)$. Bohman \cite{Boh}
proved $\alpha(F)=O(\sqrt{n\log n})$ asymptotically almost surely. Moreover $e(F)\le n\Delta(F)/2$, and
$C_4(F)\le n\Delta(F)^3/8$: a $4$-cycle is determined by a start vertex and the next three vertices along it, and it has eight
such descriptions. Proposition~\ref{p:pairing} then gives $\cm(G)\ge n/2-O(\Delta(F))-O(\Delta(F)^3/n)
=n/2-O(\sqrt n\,\log^{3/2}n)$. If $n\ge4t-1$ then $t\le(n+1)/4$, which is less than this bound for large $n$.
\end{proof}

\begin{remark}
Random pairings with these two probabilities, and the count of non-touching couples through $4$-cycles of the complement,
are used by Norin and Seymour \cite{NS} and by Yip \cite{Yip}, so Proposition~\ref{p:pairing} is a short corollary of
their counts. The application to the triangle-free process is the new step.
By Proposition~\ref{p:pairing}, every sequence of graphs with $\cm(G)=o(|V(G)|)$ and $e(\overline G)=o(|V(G)|^2)$ has
$C_4(\overline G)\ge(1/2-o(1))|V(G)|^3$.
\end{remark}


\subsection{Three remarks checked by one reader}
Each remark below was checked line by line by one reader other than its author. They are paper arguments.

\begin{remark}[inflations with two clique orders]\label{r:infl}
An \emph{inflation} of a graph $B$ replaces each vertex $x$ by a clique of order $c_x\ge0$, and joins the cliques of $x$ and
$y$ completely when $xy\in E(B)$; it is \emph{proper} if every $c_x\ge1$ and \emph{$r$-uniform} if every $c_x=r$; write
$B[K_r]$ for the $r$-uniform inflation. Costa, Luu, Wood and Yip \cite[Conjecture~1.14]{CLWY} conjecture that if $\alpha(B)=2$
and $B$ satisfies a property from a list that includes $h\ge\chi$ and $h\ge|V|/2$, then every inflation of $B$ satisfies it.
Uniform inflations keep $h\ge\chi$: $r$ copies of a clique model of $B$ give one of $B[K_r]$, and
$\chi(B[K_r])\le r\chi(B)$. Two consecutive clique orders suffice to break it:

\emph{For every integer $s\ge1$ there are arbitrarily large graphs $B$ with $\alpha(B)=2$ and $h(B)\ge\chi(B)=|V(B)|/2$ that have a
proper inflation $I$, with all cliques of order $s$ or $s+1$, such that $h(I)<|V(I)|/2\le\chi_f(I)\le\chi(I)$.} So
Conjecture~1.14 of \cite{CLWY} fails for $h\ge\chi$ and for $h\ge|V|/2$.

\emph{Projection.} For every graph $G$ and $r\ge1$, $h(G[K_r])\le r\,h_f(G)$. Indeed the projections to $G$ of the branch sets
of a $K_t$ model in $G[K_r]$ are connected and pairwise touch, and each vertex $v$ of $G$ lies in at most $r$ of them because
the branch sets are disjoint and the clique of $v$ has $r$ vertices; weight $1/r$ on each branch set gives an admissible
weighting of total $t/r$. (Fox \cite{Fox11} proves an equality of this kind; we need only this inequality.)

\emph{Construction.} Let $G$ be a graph of Theorem~\ref{t:main} with $n$ vertices and $\omega=\omega(G)\le2\cm(G)+1$. Take a
maximal family of $k$ vertex-disjoint induced paths on three vertices. The remaining vertices span a graph without induced
three-vertex paths, so each of its components is a clique, and $\alpha(G)\le2$ allows at most two components; hence
$n-2\omega\le3k\le n$. Let $a=n-2k$, add a clique $A$ of $a$ new vertices, fix a bijection $\pi$ from $A$ onto the $n-2k$ old
vertices that are not ends of the $k$ paths, and join each $x\in A$ to every old vertex except $\pi(x)$. Call the result $B$.
A new vertex $x$ has exactly one non-neighbour, $\pi(x)$, so $\alpha(B)=2$. The ends of each path and the pairs
$\{x,\pi(x)\}$ split $V(B)$ into $k+a$ independent pairs, so $\chi(B)=k+a=|V(B)|/2$. Every old vertex outside a path is
adjacent to one of its two ends (otherwise three independent vertices), and every new vertex is adjacent to both ends. So the
$k$ paths followed by the $a$ new vertices form a clique model, in fact a dominating one in the sense of \cite{KWY}, of
order $k+a$; hence $h(B)\ge\chi(B)$.

\emph{Inflation.} Let $I$ replace each old vertex by a clique of order $s+1$ and each new vertex by a clique of order $s$, and
let $N=(s+1)n+sa$. In a clique model of $I$ at most $sa$ branch sets meet the cliques of new vertices; the others form a clique
model of $G[K_{s+1}]$. By the projection inequality and Proposition~\ref{p:hf},
\[
 h(I)\le sa+(s+1)h_f(G)\le sa+\tfrac{s+1}3\bigl(n+4\cm(G)+2\bigr).
\]
Since $a\le n-\tfrac23(n-2\omega)=(n+4\omega)/3$, this gives
$N/2-h(I)\ge n/6-\tfrac{2s}3\omega-\tfrac{s+1}3(4\cm(G)+2)$, which is positive when $\beta$ is small in terms of $s$ and $n$
is large. As $\alpha(I)=2$, $\chi_f(I)\ge N/2>h(I)$.
\end{remark}

\begin{remark}[a density gap for half-order minors]\label{r:dense}
Norin and Seymour \cite{NS} proved that every graph with $\alpha(G)=2$ has a minor with $\lceil|V(G)|/2\rceil$ vertices and edge
density at least $\gamma-o(1)$, where $\gamma=0.98688\ldots$. The best constant in that theorem is less than $1$:
\emph{there are $\eps>0$ and infinitely many graphs $I$ with $\alpha(I)=2$ in which every minor with at least $|V(I)|/2$
vertices has edge density at most $1-\eps$.}

Fix a graph $B$ with $q$ vertices, $\alpha(B)=2$ and $\omega(B)+\cm(B)\le q/64$ (Theorem~\ref{t:main} gives one), and let
$I=B[K_r]$, $n=rq$. Take a clique model in $I$ with $t\ge n/2$ branch sets, and let $E$ be the number of non-adjacent pairs of
branch sets. If $u$ branch sets have at most two vertices, then $n\ge u+3(t-u)$, so $u\ge(3t-n)/2\ge n/4$. Keep each of these independently with
probability $1/(4r)$; the expected number kept is at least $q/16$. At most $r$ such branch sets meet the clique of a fixed
vertex of $B$, so the expected number of kept pairs whose projections to $B$ intersect is at most $q\binom r2/(16r^2)\le q/32$,
and the expected number of kept non-adjacent pairs is at most $E/(16r^2)$. Delete one member of each such kept pair. The
projections of the remaining kept branch sets are pairwise disjoint and pairwise adjacent sets with one or two vertices: the
one-vertex sets form a clique and the two-vertex sets form a connected matching of $B$, so at most $\omega(B)+\cm(B)\le q/64$
remain. Taking expectations, $q/16-q/32-E/(16r^2)\le q/64$, so $E\ge r^2q/4=n^2/(4q)$, and the minor misses at least
$E/\binom t2\ge1/(2q)$ of its pairs. So $\eps=1/(2q)$ works for every $r$. (OpenAI's fixed bound already gives a base $B$ for
the same argument with other constants.)
\end{remark}

\begin{remark}[strongly odd dominating models]\label{r:sod}
Kawarabayashi, Wood and Yip \cite{KWY} define a \emph{strongly odd dominating} $K_t$-model as disjoint subtrees
$T_1,\dots,T_t$ with a red/blue colouring of their vertices, such that every edge of each $T_i$ has both colours and, for
$i<j$, every vertex of $T_j$ has a neighbour of its own colour in $T_i$; $\mathrm{had}_{\rm sod}(G)$ is the largest such $t$.
They prove $\mathrm{had}_{\rm sod}(G)\ge|V(G)|/(3\alpha(G)-1)$ and give graphs with $\alpha=2$ and
$\mathrm{had}_{\rm sod}\le|V|/4+O(\sqrt{|V|\log|V|})$. For $\alpha=2$ the constant $1/4$ is the right one:
\emph{every nonempty graph with $\alpha(G)\le2$ has $\mathrm{had}_{\rm sod}(G)\ge\max\{\lceil|V(G)|/4\rceil,
\lceil\chi(G)/2\rceil\}$, and every graph has $\mathrm{had}_{\rm sod}(G)\le|V(G)|/4+\omega(G)$}; so both constants are
sharp on graphs with $\alpha=2$ and $\omega=o(|V|)$.

Lower bound: colour an induced path $abcd$ red, blue, red, blue. Its red vertices $a,c$ are non-adjacent, as are its blue
vertices $b,d$, so every other vertex has a neighbour of each colour on it ($\alpha\le2$). Remove induced four-vertex paths
one after another, $k$ in all, until the rest $H$ has none. Then $H$ is a cograph, hence perfect, so
$\omega(H)=\chi(H)\ge|V(H)|/2$. The $k$ paths followed by a maximum clique of $H$ as one-vertex red trees form a strongly odd
dominating model of order $t=k+\omega(H)\ge(4k+|V(H)|)/4=|V(G)|/4$, and $\chi(G)\le2k+\chi(H)\le2t$.

Upper bound: one-vertex trees must all have the same colour (a later one-vertex tree needs a neighbour of its own colour in
an earlier one). For each colour, the trees with exactly one vertex of that colour have these vertices pairwise adjacent
(the later one must see the earlier one). If $a_i$ trees have $i$ vertices, this gives $a_1\le\omega$ and
$a_1+2a_2+a_3\le2\omega$ (a two-vertex tree has one vertex of each colour, a three-vertex path has one colour once). Then
$|V(G)|\ge4t-3a_1-2a_2-a_3\ge4t-4\omega$.
\end{remark}

\section{A second route: the Colin de Verdi\`ere paper}\label{s:second}
A second derivation, obtained independently (by GPT-6.2 Astra, before the first route was compared with it), uses the
companion paper \cite{OAIC}. There the base graphs satisfy (G1): matchings of size $m/20$ split into classes of size at most
$am$ contain conflicts between different classes, with a fixed $a>0$ \cite[thm:base-graphs]{OAIC}. The fixed scale again comes
from a fixed endpoint cap, entering the coloured terminal cover as $\pi(S)<8/M$. The route replaces the symmetric
distribution test by one with a prescribed endpoint cap $B$, choosing $B$ before the channel size $h$. Its key point is that
the retained fraction of the symmetric phase theorem, $d_{\rm ph}=p_0d_{\rm cov}/(64R_0)$, and hence the colour threshold
$\eta=s^2/512$ with $s=10^{-6}d_{\rm ph}/64$, do not depend on the cap, while the cap enters only lower bounds on $h$ and
the size of $n$. The transfer then runs at density $\gamma$ with $B=\max\{2^{1000},128e/\gamma\}$, $a=\gamma\eta/64$ and
$k=\lceil\gamma m/4\rceil$. This route needs colours, a palette, thinning and an asymmetric test, which the Hadwiger route of
Sections~\ref{s:uniform}--\ref{s:transfer} avoids; it has no formal counterpart, since \cite{OAIL} does not cover \cite{OAIC}.
Both routes rest on the long supersaturation arguments of OpenAI.

\section{Verification state}\label{s:verify}
\paragraph{Lean.} The statements of Theorem~\ref{t:main} for rational $\beta$ (its second sentence is the same statement,
since two disjoint edges conflict exactly when they do not touch), of the upper bound in Corollary~\ref{c:third} and of
Corollary~\ref{c:conj} are proved in Lean~4.34.1 with Mathlib (commit \texttt{d13f23b}). The exported bound on the
Hadwiger number is $h<(1/3+2\beta)m+2/3$; the bound $(1/3+\tfrac43\beta)m+2/3$ of the abstract follows from
Lemma~\ref{l:mm}, which is in OpenAI's Lean library, and $\cm<\beta m$. The paper's probability estimate in
Lemma~\ref{l:transfer} is not part of the Lean statements. The
proof consists of four new modules \texttt{OAI.Combinatorics.HadwigerScaled.\{ScaledParameters, ScaledTransfer,
ScaledExponential, Main\}} on top of 288 unchanged files of OpenAI's Lean library \cite{OAIL}. Three of the new modules are
copies of OpenAI files with the cap $2^{1000}$ replaced by $2^{12Q+1}$ and the threshold $m/100$ by $m/Q$. The build of the
closure (9,216 jobs) ends with exit code~0 and has no \texttt{sorry} outside the statement file. For each of the five
theorems \texttt{exists\_graphs\_div}, \texttt{exists\_graphs\_beta}, \texttt{exists\_small\_hadwiger},
\texttt{not\_linear\_cm} and \texttt{not\_hc\_eps}, \texttt{\#print axioms} lists only \texttt{propext},
\texttt{Classical.choice} and \texttt{Quot.sound}. A statement file that imports only Mathlib states the last three theorems
with OpenAI's definitions of connected matchings and clique minors; the fully elaborated text of these definitions and
statements is byte-identical in the statement file and in the proof. Two further checks accept the proof:
\texttt{leanchecker}, which replays the Lean kernel, replayed all 292 modules, and \texttt{lean4lean}, an independent
kernel, checked an export of the five theorems with all their dependencies (39,086 declarations). A second build from an
empty folder, with the files as published and OpenAI's files fetched at the pinned commit, repeated all these checks with
the same results. The Lean statements give, for each rational $\beta>0$, infinitely many orders $m$ with a
graph satisfying $\alpha\le2$, $\cm<\beta m$ and $h<(1/3+2\beta)m+2/3$; the repository file \texttt{FIDELITY.md} compares them
with the published wording of the conjectures.

\paragraph{Paper checks.} A separate Claude agent read the second route (Section~\ref{s:second}) line by line against
\cite{OAIC} and found no error; it found independently that the Lean development of \cite{OAIH} is parametric in the
marginal cap. Two blind checks of the bare statement of Proposition~\ref{p:uniform}, given without our proof, were run.
GPT-6.2 Astra, with read access to the TeX and Lean sources of \cite{OAIH}, judged it true, listed every use of the cap, and
named as weakest step the fresh-probe deletion, whose prefactor must be controlled before $h$ is chosen (it is, by the
single-frame transpose bound $2^{d-h+1}$). Gemini 4 Neon, given only Sections~3 and 14 of \cite{OAIH}, proved
Lemma~\ref{l:transfer} in full and found no obstruction to Proposition~\ref{p:uniform}, but could not decide it from those two
sections alone. None of these checks is a referee report. The paper proofs of \cite{OAIH} and \cite{OAIC} have not been
refereed; the Lean proof does not depend on them.

\paragraph{Novelty.} Searches made on 25 November 2026, in the sources listed here and only there, found no public
statement that Linear-CM or $\mathrm{HC}_\eps$ fails,
or that $\inf h/|V|=1/3$ for $\alpha\le2$: arXiv listings and full texts of the closest papers (including Xie,
arXiv:2610.17090, and Kawarabayashi, Wood and Yip, arXiv:2611.11205, who state that the constant $1/3$ cannot be replaced by
any constant above $26/75$), the AI-mathematics tracker vibemathed.com (which lists the weak conjecture as not settled),
the full \TeX{} of 1,550 of the 1,560 works on hexagonmath.org, the Palomar registry, nilradical.ai, proofatlas.ai, Zenodo,
GitHub issues and pull requests, MathOverflow, the lists of papers that cite \cite{FGS}, \cite{Fox} and \cite{KPT}, and the
arXiv math.CO listings of 19 to 25 November 2026.

\section{Open problems}
\begin{enumerate}
\item \emph{Rate.} Our graphs have $m=2^{C_0gM_0n}$ vertices with $M_0$ depending on $\beta$. Is there one sequence with
$\cm(G)\le f(m)\,m$ and an explicit $f\to0$? Letting $M$ grow with $n$ would need $M=2^{o(N)}$ and uniform versions of several
$o(1)$ steps of \cite{OAIH}.
\item \emph{True order.} Fox's lower bound is $\cm(G)=\Omega(m^{4/5}(\log m)^{1/5})$; is $\min\cm$ over $m$-vertex graphs with
$\alpha\le2$ of order $m^{1-c}$ for some $c>0$? Equivalently, what is the order of $\min h(G)-m/3$?
\item \emph{Explicit graphs.} Are there explicit or algebraic graphs with $\alpha\le2$ and $\cm=o(m)$, of reasonable size?
\end{enumerate}

\begin{thebibliography}{99}
\bibitem{BK} J.~Balogh and A.~Kostochka, Large minors in graphs with given independence number, \emph{Discrete Math.}
\textbf{311} (2011), 2203--2215.
\bibitem{Bon} GitHub user \texttt{EdouardBonnet}, \texttt{hadwiger-counterexample}: a Lean~4 port of OpenAI's
formalization of their counterexample to Hadwiger's conjecture ($100\,\cm(G)<m$), GitHub repository
\texttt{EdouardBonnet/hadwiger-counterexample}, kernel replay dated 11 October 2026.
\bibitem{BKT} T.~B\"ohme, A.~Kostochka and A.~Thomason, Minors in graphs with high chromatic number, \emph{Combin.
Probab. Comput.} \textbf{20} (2011), 513--518.
\bibitem{CD} R.~Chen and Z.~Deng, Connected matchings in graphs with independence number two, \emph{Graphs Combin.}
\textbf{41} (2025), Article 107; arXiv:2409.05920.
\bibitem{CF} M.~Chudnovsky and A.~O.~Fradkin, An approximate version of Hadwiger's conjecture for claw-free graphs,
\emph{J. Graph Theory} \textbf{63} (2010), 259--278.
\bibitem{CS} M.~Chudnovsky and P.~Seymour, Packing seagulls, \emph{Combinatorica} \textbf{32} (2012), 251--282.
\bibitem{Cam} S.~Cambie, Hadwiger's conjecture implies a conjecture of F\"uredi--Gy\'arf\'as--Simonyi, arXiv:2108.10303 (2021).
\bibitem{CFS} D.~Conlon, J.~Fox and B.~Sudakov, Short proofs of some extremal results III, arXiv:1910.08661 (2019).
\bibitem{CLWY} J.~Costa, E.~Luu, D.~R.~Wood and J.~H.~Yip, Verifying Hadwiger's conjecture for examples of graphs with
$\alpha(G)=2$, arXiv:2512.17114 (2025).
\bibitem{DM} P.~Duchet and H.~Meyniel, On Hadwiger's number and the stability number, in \emph{Graph Theory (Cambridge, 1981)},
North-Holland Math.\ Stud.\ \textbf{62} (1982), 71--73.
\bibitem{Fox} J.~Fox, Complete minors and independence number, \emph{SIAM J. Discrete Math.} \textbf{24} (2010), 1313--1321.
\bibitem{FGS} Z.~F\"uredi, A.~Gy\'arf\'as and G.~Simonyi, Connected matchings and Hadwiger's conjecture,
\emph{Combin. Probab. Comput.} \textbf{14} (2005), 435--438.
\bibitem{Had} H.~Hadwiger, \"Uber eine Klassifikation der Streckenkomplexe, \emph{Vierteljschr. Naturforsch. Ges. Z\"urich}
\textbf{88} (1943), 133--142.
\bibitem{KPT} K.~Kawarabayashi, M.~D.~Plummer and B.~Toft, Improvements of the theorem of Duchet and Meyniel on Hadwiger's
conjecture, \emph{J. Combin. Theory Ser. B} \textbf{95} (2005), 152--167.
\bibitem{KWY} K.~Kawarabayashi, D.~R.~Wood and J.~H.~Yip, Independence number and odd dominating models,
arXiv:2611.11205v2 (10 November 2026).
\bibitem{OAIH} OpenAI, A counterexample to Hadwiger's conjecture, preprint, 23 September 2026, repository
\texttt{openai/math}, \texttt{preprints/A-counterexample-to-Hadwigers-conjecture-September-23-2026}.
\bibitem{OAIC} OpenAI, A counterexample to the Colin de Verdi\`ere chromatic conjecture, preprint, 23 September 2026, same repository.
\bibitem{OAIL} OpenAI, Lean library of the \texttt{openai/math} repository, commit \texttt{3344da59991ed94f9ba711217b29c1d760badfa6}, modules
\texttt{OAI.Combinatorics.HadwigerCounterexample}; scope file \texttt{lean/docs/157.md}.
\bibitem{PST} M.~D.~Plummer, M.~Stiebitz and B.~Toft, On a special case of Hadwiger's conjecture, \emph{Discuss. Math. Graph
Theory} \textbf{23} (2003), 333--363.
\bibitem{Sey} P.~Seymour, Hadwiger's conjecture, in \emph{Open Problems in Mathematics} (J.~F.~Nash Jr.\ and M.~Th.~Rassias,
eds.), Springer, 2016, 417--437.
\bibitem{Yip} J.~H.~Yip, Dense matchings of linear size in graphs with independence number 2, arXiv:2512.01401 (2025).
\bibitem{Boh} T.~Bohman, The triangle-free process, \emph{Adv. Math.} \textbf{221} (2009), 1653--1677.
\bibitem{Fox11} J.~Fox, Constructing dense graphs with sublinear Hadwiger number, arXiv:1108.4953 (2011).
\bibitem{HW} D.~J.~Harvey and D.~R.~Wood, Parameters tied to treewidth, \emph{J. Graph Theory} (2017),
  doi:10.1002/jgt.22030; arXiv:1312.3401v2.
\bibitem{NS} S.~Norin and P.~Seymour, Dense minors of graphs with independence number two, \emph{J. Combin. Theory
  Ser. B} \textbf{176} (2026), 101--110.
\end{thebibliography}

\end{document}
